\expandafter\def\csname HistoricalNoteText15\endcsname{Dass sie allemal einen ganz bestimmten Lauf hat, beruhet darauf, dass sie einen Theil der äussern
Abgrenzung einer Fläche, für welche $T$ ein bestimmtes Zeichen hat, ausmachen soll: ich habe das positive
Zeichen gewählt, was an sich ganz willkürlich ist. So verstanden setzt sich z. B. die in $1$ eintretende Linie
durch $11$ nach $10$ fort: als Theil der Grenzlinie einer Fläche, worin $T$ negativ ist, würde die Linie $1.11$
nach $2$ fortgesetzt werden müssen. Spricht man hingegen nur von einer Linie worin $T=0$ ist, ohne sie als
Theil der Begrenzung einer bestimmten Fläche zu betrachten, so würde eher $11.9$ als natürliche Fortsetzung
von $1.11$ gelten können. Der hier gewählte Gesichtspunkt unterscheidet mein gegenwärtiges Verfahren von
dem von $1799$, und trägt wesentlich zur Vereinfachung der Beweisführung bei.}
